\documentclass[11pt]{article}
\usepackage{../homework}

\profname{Dr. David Singer}
\classcode{MATH 465 Differential Geometry}
\assignnum{Assignment 05}
\duedate{5 October 2026}

\begin{document}
All problems from Do Carmo, \textit{Differential Geometry of Curves and Surfaces} (2e).

\begin{problem}{3.2.5}
  Show that the mean curvature \( H \) at \( p \in S \) is given by
  \begin{equation}
  \label{eq:1}
  H = \frac{1}{\pi} \int\limits_0^{\pi} k_n(\theta) d \theta,
\end{equation}
where \( k_n(\theta) \) is the normal curvature at \( p \) along a direction making an angle \(\theta\) with a fixed direction.
\end{problem}
\begin{proof}
  Let \( L: T_pS \rightarrow T_pS \) be the Weingarten map (equivalently, \( -dN_p \) for \( N \) the Gauss map), which, by the spectral theorem for finite dimensional operators, has an orthonormal set of eigenvectors \( e_1, e_2 \), the principal directions of curvature of \( S \) at \( p \), with associated principal curvatures \( k_1 \) and \( k_2 \).

  If \(\alpha: I \rightarrow S\) is any curve on the surface with \(\alpha(0)=p\) and \( \alpha^{\prime}(0) = V \), then \( V \in T_pS \) and we can write \( V = \cos\theta e_1 + \sin\theta e_2 \).
  The normal curvature along \( V \) is \( k_n = D_VV \cdot \nu = \langle L(V), V \rangle \).
  \begin{equation}
  \label{eq:9}
  \begin{aligned}
    \langle L(V), V \rangle &= \langle L(\cos\theta e_1 + \sin\theta e_2), \cos\theta e_1 + \sin\theta e_2 \rangle \\
                &= \langle k_1\cos\theta e_1 + k_2\sin\theta e_2, \cos\theta e_1 + \sin\theta e_2 \rangle \\
    &= k_1\cos^2\theta + k_2 \sin^2\theta.
  \end{aligned}
\end{equation}

Therefore, we can write \cref{eq:1} as
\begin{equation}
\label{eq:10}
\begin{aligned}
  \frac{1}{\pi} \int\limits_0^{\pi} (k_1\cos^2\theta + k_2\sin^2\theta) d\theta &= \frac{1}{\pi} \left[ \int\limits_0^{\pi} k_1 \frac{1 + \cos2\theta}{2}d\theta + \int\limits_0^{\pi} k_2 \frac{1-\cos2\theta}{2}d\theta \right] \\
  &= \frac{1}{\pi} \left( k_1( \frac{\pi}{2} ) + k_2( \frac{\pi}{2} ) \right) = \frac{k_{1} + k_2}{2}.
\end{aligned}
\end{equation}
\end{proof}

\begin{problem}{3.2.8}
  Describe the region of the unit sphere covered by the image of the Gauss map of the following surfaces:
  \begin{enumerate}[label=(\alph*)]
    \item Paraboloid of revolution \( z = x^2 + y^2 \),
          \begin{answer}
            The normal vector at the origin is mapped to \( (0,0,-1) \subset S^2 \).
            Traversing upwards in \( z \), we see that the range of the Gauss map is exactly the open subset making up the lower hemisphere of \( S^2 \), that is,
            \begin{equation}
            \label{eq:11}
            N(S) = \left\{ (x,y,z) \subset S^2 \mid z < 0 \right\}.
            \end{equation}
          \end{answer}
    \item Hyperboloid of revolution \( x^2 + y^2 - z^2 = 1 \),
          \item Catenoid \( x^2 + y^2 = \cosh^2 z \).
  \end{enumerate}
\end{problem}

\begin{problem}{3.2.13 (Theorem of Beltrami-Enneper)}
  Prove that the absolute value of the torsion \(\tau\) at a point of an asymptotic curve, whose curvature is nowhere zero, is given by
  \begin{equation}
  \label{eq:2}
  | \tau | = \sqrt{- K},
\end{equation}
where \( K \) is the Gaussian curvature of the surface at a given point.
\end{problem}

\begin{problem}{3.3.3}
  Determine the asymptotic curves of the catenoid
  \begin{equation}
  \label{eq:3}
  \mathbf{x}(u,v) = (\cosh v \cos u, \cosh v \sin u, v).
  \end{equation}
\end{problem}

\begin{problem}{3.3.5}
  Consider the parametrized surface (Enneper's surface)
  \begin{equation}
  \label{eq:4}
  \mathbf{x}(u,v) = (u - \frac{u^3}{3} + uv^2, v - \frac{v^3}{3} + vu^2, u^2 - v^2),
\end{equation}
and show that
\begin{enumerate}[label=(\alph*)]
  \item The coefficients of the first fundamental form are
        \begin{equation}
        \label{eq:5}
        E = G = (1 + u^2 + v^2)^2, \quad F = 0.
      \end{equation}
  \item The coefficients of the second fundamental form are
        \begin{equation}
        \label{eq:6}
        e = 2, \quad g = -2, \quad f = 0.
      \end{equation}
  \item The principal curvatures are
        \begin{equation}
        \label{eq:7}
        k_1 = \frac{2}{(1 + u^2 + v^2)^2}, \quad k_2 = - \frac{2}{(1 + u^2 + v^2)^2}.
      \end{equation}
  \item The lines of curvature are the coordinate curves.
  \item The asymptotic curves are \( u + v = \text{const.} \), \( u - v = \text{const.} \)
\end{enumerate}
\end{problem}
\end{document}
